\subsection{ORBITS}

PROPOSITION 14. Let \(\mathbf{G}\) be a Lie group, \(\mathbf{X}\) an analytic manifold and \((g,x)\mapsto gx\) a law of analytic left operation of \(\mathbf{G}\) on \(\mathbf{X}\) . Let \(x\in \mathbf{X}\) . Suppose that the corresponding orbital mapping \(\rho (x)\) is a subimmersion (which is always the case if \(\mathbf{K}\) is of characteristic 0 and \(\mathbf{X}\) is finite-dimensional (Corollary to Proposition 9)). Let \(\mathbf{G}_{\mathbf{x}}\) be the stabilizer of \(x\) in \(\mathbf{G}\) .

\begin{enumerate}
\def\labelenumi{(\roman{enumi})}
\item
  \(\mathrm{G}_x\) is a Lie subgroup and \(\mathrm{T}_e(\mathrm{G}_x) = \mathrm{Ker}\, \mathrm{T}_e(\rho(x))\) .
\item
  The canonical mapping \(i_x\) of the homogeneous space \(G / G_x\) into \(X\) is an immersion with image \(Gx\) .
\item
  If further the orbit \(\mathbf{G}x\) is locally closed and the topology on \(\mathbf{G}\) admits a countable base, then \(\mathbf{G}x\) is a submanifold of \(\mathbf{X}\) , \(i_x\) is an isomorphism of the manifold \(\mathrm{G} / \mathrm{G}_x\) onto the manifold \(\mathbf{G}x\) and \(\mathrm{T}_x(\mathbf{G}x) = \operatorname{Im} \mathrm{T}_e(\rho(x))\) .
\end{enumerate}

The inverse image of \(x\) under \(\rho(x)\) is \(G_x\) . As \(\rho(x)\) is a subimmersion, \(G_x\) is a submanifold and, for all \(g \in G\) , the tangent space \(J\) to \(gG_x = \rho(x)^{-1}(gx)\) at \(g\) is \(\operatorname{Ker} T_g(\rho(x))\) (Differentiable and Analytic Manifolds, R, 5.10.5), whence (i). Let \(\pi: G \to G/G_x\) be the canonical mapping. Then \(i_x \circ \pi = \rho(x)\) . As \(G/G_x\) is a quotient manifold of \(G\) , this equality proves that \(i_x\) is analytic. Further, the kernels of \(T_g(\rho(x))\) and \(T_g(\pi)\) are both equal to \(J\) . Hence \(T_{\pi(g)}(i_x)\) is injective. The image of \(T_{\pi(g)}(i_x)\) is equal to the image of \(T_g(\rho(x))\) and hence admits a topological supplement. This proves (ii).

Suppose that \(Gx\) is locally closed. Every point of \(Gx\) then has a neighbourhood in \(Gx\) which is homeomorphic to a closed subspace of a complete metric space and hence is a Baire space. Hence \(Gx\) is a Baire space (General Topology, Chapter IX, § 5, Proposition 4). If \(G\) has a countable base, \(i_x\) is therefore a homeomorphism of \(G / G_x\) onto \(Gx\) (General Topology, Chapter IX, § 5). Then by (ii) and Differentiable and Analytic Manifolds, R, 5.8.3, \(i_x\) is an isomorphism of the manifolds \(G / G_x\) onto the manifold \(Gx\) and

\[
\mathrm{T} _ {x} (\mathrm{G} x) = \operatorname{Im} \mathrm{T} _ {\pi (e)} (i _ {x}) = \operatorname{Im} \mathrm{T} _ {e} (\rho (x)).
\]

Remark. Let G be a finite-dimensional Lie group, X a manifold of class \(C^{r}\) and \((g, x) \mapsto gx\) a law of left operation of class \(C^{r}\) of G on X. Then Proposition 14 remains valid. The only point which needs a different proof is the fact that \(G_{x}\) is a Lie subgroup. But if \(r \neq \omega\) , K = R; as clearly \(G_{x}\) is closed, \(G_{x}\) is a Lie subgroup by §8, Theorem 2.

COROLLARY. Let G be a Lie group whose topology admits a countable base and X a

non-empty finite-dimensional analytic manifold with a law of analytic left operation of G on X. Suppose that G operates transitively on X and that K is of characteristic 0. Then X is a Lie homogeneous space for G.

Let \(x \in X\) . The orbit of \(x\) , equal to \(X\) , is closed and we can therefore apply Proposition 14 (iii).

